% Parte do material do minicurso "Escrita de Artigos Científicos com LaTeX" (Conecta Hub)
% Exemplo 03 – Matemática

\documentclass[12pt, a4paper]{article}
\usepackage[utf8]{inputenc}
\usepackage[T1]{fontenc}
\usepackage[brazil]{babel}
\usepackage{amsmath, amssymb}            % pacotes essenciais para matemática

\title{Matemática no LaTeX}
\author{Seu Nome}
\date{}

\begin{document}

\maketitle

\section{Equação inline}
A equação do segundo grau é $ax^2 + bx + c = 0$, com $a \neq 0$.

\section{Equação em destaque}
Equação numerada:
\begin{equation}
    x = \frac{-b \pm \sqrt{b^2 - 4ac}}{2a}
    \label{eq:bhaskara}
\end{equation}

Veja a equação~\eqref{eq:bhaskara} para a fórmula de Bhaskara.

Equação sem numeração (\texttt{equation*}):
\[
    \int_0^1 x^2 \, dx = \frac{1}{3}
\]

\section{Várias linhas alinhadas}
\begin{align}
    (a + b)^2 &= a^2 + 2ab + b^2 \\
    (a - b)^2 &= a^2 - 2ab + b^2
\end{align}

\section{Símbolos úteis}
Letras gregas: $\alpha, \beta, \gamma, \Delta, \Omega$.

Somatório e produtório: $\sum_{i=1}^{n} i = \frac{n(n+1)}{2}$ e $\prod_{i=1}^{n} i = n!$.

Conjuntos: $x \in A$, $A \subseteq B$, $\mathbb{R}$, $\mathbb{N}$.

Matrizes:
\[
    A = \begin{pmatrix} 1 & 2 \\ 3 & 4 \end{pmatrix}
\]

\end{document}
